\documentclass[../../../include/open-logic-section]{subfiles}

\begin{document}

\olfileid{sth}{cardinals}{cp}
\olsection{اصلِ کانتور}

به \olref[ordinals][vn]{sec} بازگردید. آنجا دربارهٔ مجموعه‌های
خوش‌ترتیب سخن می‌گفتیم و پیشنهاد کردیم خوب است اشیایی داشته باشیم
که جانشینِ خوش‌ترتیب‌ها شوند. با در نظر داشتن این نکته، اعداد ترتیبی
را معرفی کردیم و سپس در
\olref[ordinals][ordtype]{ordtypesworklikeyouwant} نشان دادیم که
درست به شیوهٔ مطلوب رفتار می‌کنند؛ یعنی:
\[
	\ordtype{A, <} = \ordtype{B, \lessdot}
	\text{ اگر و تنها اگر } \tuple{A, <} \isomorphic \tuple{B, \lessdot}.
\]
باز هم عقب‌تر بروید، به \olref[sfr][siz][equ]{sec}. آنجا با رویکردی
ساده‌انگارانه مفهوم «اندازه» یک مجموعه را معرفی کردیم. به‌طور مشخص،
گفتیم دو مجموعه هم‌شمارند، یعنی $\cardeq{A}{B}$، درست در صورتی که
!!a{bijection} مانند $f \colon A \to B$ وجود داشته باشد. این مفهوم
ذاتاً {ساده‌تر} از مفهوم خوش‌ترتیبی است: تنها به !!{bijection}ها
علاقه‌مندیم، نه به این‌که آیا !!{bijection}ها (مانند یکریختی‌های
ترتیبی) «ساختاری را حفظ می‌کنند» یا نه.

از همهٔ این مطالب اندیشه‌ای بدیهی زاده می‌شود. همان‌طور که برای
سنجش خوش‌ترتیب‌ها اشیای معینی، یعنی \emph{اعداد ترتیبی}، را معرفی
کردیم، می‌توانیم برای سنجش اندازه اشیای معینی، یعنی
\emph{کاردینال‌ها}، را معرفی کنیم. هدف این فصل همین است.

پیش از آن‌که بگوییم این کاردینال‌ها چه خواهند بود، باید اصلی را
مقرر کنیم که باید برآورده سازند. اگر $\card{X}$ را برای کاردینالیتهٔ
مجموعهٔ $X$ بنویسیم، می‌خواهیم:
\[
	\card{A} = \card{B} \text{ اگر و تنها اگر } \cardeq{A}{B}.
\]
این را \emph{اصل کانتور} می‌نامیم، زیرا احتمالاً کانتور نخستین کسی
بود که آن را کاملاً روشن در نظر داشت. (در
\olref[hp]{sec} دربارهٔ رابطهٔ آن با \emph{اصل هیوم} بیشتر سخن
خواهیم گفت.) پس هدف ما آن است که برای هر $X$، مقدار $\card{X}$ را
چنان تعریف کنیم که اصل کانتور را به دست دهد.

\end{document}
